\chapter{Basic Metabolic Panel (BMP)}

\section*{What Is This Test?}
A basic metabolic panel (BMP) is a group of 8 tests performed on a blood sample that provides important information about your body's chemical balance and kidney function. It includes: glucose (blood sugar), calcium, sodium, potassium, chloride, carbon dioxide (bicarbonate), blood urea nitrogen (BUN), and creatinine. The BMP evaluates your blood sugar, electrolyte balance, and kidney function. It is commonly ordered as a general health screening test, to monitor chronic conditions such as diabetes and kidney disease, and to evaluate symptoms such as nausea, fatigue, or confusion.

\section*{Alternative Names}
\begin{itemize}
  \item BMP
  \item Basic Metabolic Panel
  \item Chemistry Panel
  \item Chem 8
  \item Electrolyte Panel with Glucose
\end{itemize}
\section*{Principle}
The BMP is performed on automated clinical chemistry analyzers using photometric (colorimetric) and potentiometric methods. Glucose is measured by enzymatic methods (glucose oxidase or hexokinase). Electrolytes (sodium, potassium, chloride) are measured by ion-selective electrode potentiometry. BUN and creatinine are measured enzymatically. Calcium is measured by colorimetric methods or ion-selective electrodes. Carbon dioxide is measured by enzymatic or conductimetric methods. Results are typically available within 1 to 2 hours.

\section*{History}
The BMP evolved from the multi-test chemistry panels developed in the 1960s-1970s with automated analyzers. The College of American Pathologists standardized the BMP as an 8-test panel. It is a subset of the comprehensive metabolic panel (CMP, 14 tests) and is ordered more frequently because of its lower cost and focus on electrolytes and kidney function.

\section*{Why This Test Is Ordered}
\begin{itemize}
  \item Routine health screening and annual physical examination: the BMP is one of the most frequently ordered panels in medicine
  \item Evaluation of electrolyte and acid-base disturbances: nausea, vomiting, diarrhea, confusion, weakness, fatigue, or abnormal heart rhythm
  \item Assessment of kidney function in patients with hypertension, diabetes, heart failure, or known renal disease
  \item Evaluation of hyperglycemia or hypoglycemia and monitoring of glucose control
  \item Preoperative assessment as part of the surgical workup
  \item Monitoring of medications that affect electrolytes or kidney function: diuretics, ACE inhibitors, ARBs, aminoglycosides, and chemotherapy
  \item Emergency department evaluation of dehydration, intoxication, or altered mental status
  \item Monitoring of patients with chronic diseases such as heart failure, cirrhosis, and chronic kidney disease
  \item Screening for chronic kidney disease in at-risk populations (diabetes, hypertension, family history)
\end{itemize}

\section*{Primary vs Secondary Test}
This is a primary screening test for electrolyte and kidney disorders.

\section*{How This Test Is Performed}
For most patients, a health care professional will take a blood sample from a vein in your arm using a small needle. After the needle is inserted, a small amount of blood will be collected into a test tube or vial. You may feel a little sting when the needle goes in or out. This usually takes less than five minutes. For certain tests, a urine sample, stool sample, or other specimen may be required.

\section*{What Preparation Is Needed}
Preparation requirements vary by test. Some tests require fasting for 8 to 12 hours. Others have no special preparation. Your healthcare provider or laboratory will inform you of any specific preparation needed for this test.

\section*{Contraindications and Precautions}
\begin{itemize}
  \item No absolute contraindications. Fasting affects glucose. Dehydration elevates BUN, creatinine, and electrolytes. Specimen hemolysis falsely elevates potassium.
\end{itemize}
\section*{Risks}
Minimal risk. Slight pain or bruising at the venipuncture site.

\section*{What Happens After the Test}
Results are typically available within 1 to 2 hours. Abnormal results may prompt additional testing (lipid panel, liver function tests, urinalysis) or treatment adjustments for chronic conditions.

\section*{Understanding Results}

\subsection*{Below Normal}
\begin{itemize}
  \item Hypoglycemia: Fasting, insulinoma, liver disease.
  \item Hypokalemia: Diuretics, GI losses.
  \item Hyponatremia: SIADH, heart failure.
  \item Hypocalcemia: Hypoparathyroidism, vitamin D deficiency.
  \item Low bicarbonate: Metabolic acidosis.
\end{itemize}

\subsection*{Above Normal}
\begin{itemize}
  \item Elevated glucose: Diabetes, stress hyperglycemia.
  \item Elevated BUN and creatinine: Impaired kidney function.
  \item Hyperkalemia: Kidney failure, cell lysis, medications (ACE inhibitors, potassium-sparing diuretics).
  \item Hypernatremia: Dehydration.
  \item Elevated calcium: Hyperparathyroidism, malignancy.
\end{itemize}

\section*{Normal Results}
\begin{itemize}
  \item \textbf{Glucose (fasting):} 70--100 mg/dL
  \item \textbf{BUN:} 7--20 mg/dL
  \item \textbf{Creatinine:} 0.7--1.3 mg/dL
  \item \textbf{Sodium:} 136--145 mEq/L
  \item \textbf{Potassium:} 3.5--5.0 mEq/L
  \item \textbf{Chloride:} 98--106 mEq/L
  \item \textbf{Carbon dioxide:} 23--29 mEq/L
  \item \textbf{Calcium:} 8.5--10.5 mg/dL
\end{itemize}

\section*{Follow-up Tests}
\begin{itemize}
  \item \textbf{Hemoglobin A1c:} If glucose is elevated, to assess long-term glycemic control
  \item \textbf{Urinalysis:} If kidney function is abnormal, to evaluate for proteinuria and hematuria
  \item \textbf{Estimated GFR (eGFR):} Calculated from creatinine, age, and sex --- provides a more accurate assessment of kidney function \cite{levey2009new}
  \item \textbf{Comprehensive Metabolic Panel:} If additional liver function testing is warranted
  \item \textbf{Lipid panel:} Often ordered alongside BMP for cardiovascular risk assessment
\end{itemize}


\section*{References}
\bibliographystyle{unsrtnat}
\bibliography{references}

\clearpage
\section*{Regulatory Status and Authorization Date}
This chapter describes a test category, not a specific commercial product. FDA, CDSCO, EU IVDR, and MHRA approval or registration dates therefore cannot be assigned without the assay manufacturer, product name, instrument, and intended use. For laboratory-developed tests, an individual agency approval date may not exist. See the Preface for the interpretation of regulatory status.

\clearpage
